% Page 065

\seite{65}

\margmark{25. März.}%
Heute fand eine Revision der Schule durch den Herrn\ob
Kreisschulrat Mahrbach statt.

\pend
\begin{verse}
Die politische Lage wird immer spannender. Täglich
\end{verse}
\pstart
werden Deutsche von der Besatzungs-Behörde ausgewiesen,\ob
so z.\ B. in unserer Kreisstadt Bitburg die ganzen Spitzen\ob
der leitenden Behörden, an welchen der Landrat Dasbach.\ob
Die preußische Eisenbahn hat den Verkehr eingestellt.\ob
Die Folgen sind unabsehbar.

\margmark{26. März.}%
Zur Schulentlassung kamen 3 Mädchen. Neu aufge\ohb
nommen wurden ebenfalls 3 Mädchen.\ob
Die Osterferien dauern vom 28.\ März bis einschließlich\ob
9.\ April.

\margmark{28. April.}%
~~Die Natur scheint dieses Jahr Klunders wirken zu\ob
wollen.~~\ob
Am heutigen Morgen verließen zwei Sefferweicher\ob
Junggesellen namens Fischbach, Nikolaus und Müller\ob
Nikolaus ihre Heimat um in die Neue Welt nach\ob
Nordamerika zu wandern. Am 2.\ Mai werden sie\ob
Rotterdam mit dem Schiffe                    verlassen\ob
die Auswanderungszeit öffnet.

\pend
\begin{verse}
Der Gesangverein \enquote{Concordia} veranstaltete diesen den
\end{verse}
\pstart
Scheidenden eine Abschiedsfeier. Mögen sie in der Neuen\ob
Welt ihr erhofftes Glück finden!

\margmark{19. Mai.}%
Die Pfingstferien dauern vom 19.\ bis einschließ\ohb
lich 28.\ Mai. -- Politisch schärft sich die Lage immer\ob
mehr zu. Die Ausweisungen aus den besetzten Gebieten\ob
werden immer zahlreicher; der Markkurs wird immer\ob
niedriger. Die Preise steigen ins Unermeßliche. Der\ob
passive Widerstand des Volkes dauert fort.
\pend
